\IfFileExists{stacks-project.cls}{%
\documentclass{stacks-project}
}{%
\documentclass{amsart}
}

% For dealing with references we use the comment environment
\usepackage{verbatim}
\newenvironment{reference}{\comment}{\endcomment}
%\newenvironment{reference}{}{}
\newenvironment{slogan}{\comment}{\endcomment}
\newenvironment{history}{\comment}{\endcomment}

% For commutative diagrams we use Xy-pic
\usepackage[all]{xy}

% We use 2cell for 2-commutative diagrams.
\xyoption{2cell}
\UseAllTwocells

% We use multicol for the list of chapters between chapters
\usepackage{multicol}

% This is generally recommended for better output
\usepackage{lmodern}
\usepackage[T1]{fontenc}

% For cross-file-references
\usepackage{xr-hyper}

% Package for hypertext links:
\usepackage{hyperref}

% For any local file, say "hello.tex" you want to link to please
% use \externaldocument[hello-]{hello}
\externaldocument[introduction-]{introduction}
\externaldocument[conventions-]{conventions}
\externaldocument[sets-]{sets}
\externaldocument[categories-]{categories}
\externaldocument[topology-]{topology}
\externaldocument[sheaves-]{sheaves}
\externaldocument[sites-]{sites}
\externaldocument[stacks-]{stacks}
\externaldocument[fields-]{fields}
\externaldocument[algebra-]{algebra}
\externaldocument[brauer-]{brauer}
\externaldocument[homology-]{homology}
\externaldocument[derived-]{derived}
\externaldocument[simplicial-]{simplicial}
\externaldocument[more-algebra-]{more-algebra}
\externaldocument[smoothing-]{smoothing}
\externaldocument[modules-]{modules}
\externaldocument[sites-modules-]{sites-modules}
\externaldocument[injectives-]{injectives}
\externaldocument[cohomology-]{cohomology}
\externaldocument[sites-cohomology-]{sites-cohomology}
\externaldocument[dga-]{dga}
\externaldocument[dpa-]{dpa}
\externaldocument[sdga-]{sdga}
\externaldocument[hypercovering-]{hypercovering}
\externaldocument[schemes-]{schemes}
\externaldocument[constructions-]{constructions}
\externaldocument[properties-]{properties}
\externaldocument[morphisms-]{morphisms}
\externaldocument[coherent-]{coherent}
\externaldocument[divisors-]{divisors}
\externaldocument[limits-]{limits}
\externaldocument[varieties-]{varieties}
\externaldocument[topologies-]{topologies}
\externaldocument[descent-]{descent}
\externaldocument[perfect-]{perfect}
\externaldocument[more-morphisms-]{more-morphisms}
\externaldocument[flat-]{flat}
\externaldocument[groupoids-]{groupoids}
\externaldocument[more-groupoids-]{more-groupoids}
\externaldocument[etale-]{etale}
\externaldocument[chow-]{chow}
\externaldocument[intersection-]{intersection}
\externaldocument[pic-]{pic}
\externaldocument[weil-]{weil}
\externaldocument[adequate-]{adequate}
\externaldocument[dualizing-]{dualizing}
\externaldocument[duality-]{duality}
\externaldocument[discriminant-]{discriminant}
\externaldocument[derham-]{derham}
\externaldocument[local-cohomology-]{local-cohomology}
\externaldocument[algebraization-]{algebraization}
\externaldocument[curves-]{curves}
\externaldocument[resolve-]{resolve}
\externaldocument[models-]{models}
\externaldocument[functors-]{functors}
\externaldocument[equiv-]{equiv}
\externaldocument[pione-]{pione}
\externaldocument[etale-cohomology-]{etale-cohomology}
\externaldocument[proetale-]{proetale}
\externaldocument[relative-cycles-]{relative-cycles}
\externaldocument[more-etale-]{more-etale}
\externaldocument[trace-]{trace}
\externaldocument[crystalline-]{crystalline}
\externaldocument[spaces-]{spaces}
\externaldocument[spaces-properties-]{spaces-properties}
\externaldocument[spaces-morphisms-]{spaces-morphisms}
\externaldocument[decent-spaces-]{decent-spaces}
\externaldocument[spaces-cohomology-]{spaces-cohomology}
\externaldocument[spaces-limits-]{spaces-limits}
\externaldocument[spaces-divisors-]{spaces-divisors}
\externaldocument[spaces-over-fields-]{spaces-over-fields}
\externaldocument[spaces-topologies-]{spaces-topologies}
\externaldocument[spaces-descent-]{spaces-descent}
\externaldocument[spaces-perfect-]{spaces-perfect}
\externaldocument[spaces-more-morphisms-]{spaces-more-morphisms}
\externaldocument[spaces-flat-]{spaces-flat}
\externaldocument[spaces-groupoids-]{spaces-groupoids}
\externaldocument[spaces-more-groupoids-]{spaces-more-groupoids}
\externaldocument[bootstrap-]{bootstrap}
\externaldocument[spaces-pushouts-]{spaces-pushouts}
\externaldocument[spaces-chow-]{spaces-chow}
\externaldocument[groupoids-quotients-]{groupoids-quotients}
\externaldocument[spaces-more-cohomology-]{spaces-more-cohomology}
\externaldocument[spaces-simplicial-]{spaces-simplicial}
\externaldocument[spaces-duality-]{spaces-duality}
\externaldocument[formal-spaces-]{formal-spaces}
\externaldocument[restricted-]{restricted}
\externaldocument[spaces-resolve-]{spaces-resolve}
\externaldocument[formal-defos-]{formal-defos}
\externaldocument[defos-]{defos}
\externaldocument[cotangent-]{cotangent}
\externaldocument[examples-defos-]{examples-defos}
\externaldocument[algebraic-]{algebraic}
\externaldocument[examples-stacks-]{examples-stacks}
\externaldocument[stacks-sheaves-]{stacks-sheaves}
\externaldocument[criteria-]{criteria}
\externaldocument[artin-]{artin}
\externaldocument[quot-]{quot}
\externaldocument[stacks-properties-]{stacks-properties}
\externaldocument[stacks-morphisms-]{stacks-morphisms}
\externaldocument[stacks-limits-]{stacks-limits}
\externaldocument[stacks-cohomology-]{stacks-cohomology}
\externaldocument[stacks-perfect-]{stacks-perfect}
\externaldocument[stacks-introduction-]{stacks-introduction}
\externaldocument[stacks-more-morphisms-]{stacks-more-morphisms}
\externaldocument[stacks-geometry-]{stacks-geometry}
\externaldocument[moduli-]{moduli}
\externaldocument[moduli-curves-]{moduli-curves}
\externaldocument[examples-]{examples}
\externaldocument[exercises-]{exercises}
\externaldocument[guide-]{guide}
\externaldocument[desirables-]{desirables}
\externaldocument[coding-]{coding}
\externaldocument[obsolete-]{obsolete}
\externaldocument[fdl-]{fdl}
\externaldocument[index-]{index}

% Theorem environments.
%
\theoremstyle{plain}
\newtheorem{theorem}[subsection]{Theorem}
\newtheorem{proposition}[subsection]{Proposition}
\newtheorem{lemma}[subsection]{Lemma}

\theoremstyle{definition}
\newtheorem{definition}[subsection]{Definition}
\newtheorem{example}[subsection]{Example}
\newtheorem{exercise}[subsection]{Exercise}
\newtheorem{situation}[subsection]{Situation}

\theoremstyle{remark}
\newtheorem{remark}[subsection]{Remark}
\newtheorem{remarks}[subsection]{Remarks}

\numberwithin{equation}{subsection}

% Macros
%
\def\lim{\mathop{\mathrm{lim}}\nolimits}
\def\colim{\mathop{\mathrm{colim}}\nolimits}
\def\Spec{\mathop{\mathrm{Spec}}}
\def\Hom{\mathop{\mathrm{Hom}}\nolimits}
\def\Ext{\mathop{\mathrm{Ext}}\nolimits}
\def\SheafHom{\mathop{\mathcal{H}\!\mathit{om}}\nolimits}
\def\SheafExt{\mathop{\mathcal{E}\!\mathit{xt}}\nolimits}
\def\Sch{\mathit{Sch}}
\def\Mor{\mathop{\mathrm{Mor}}\nolimits}
\def\Ob{\mathop{\mathrm{Ob}}\nolimits}
\def\Sh{\mathop{\mathit{Sh}}\nolimits}
\def\NL{\mathop{N\!L}\nolimits}
\def\CH{\mathop{\mathrm{CH}}\nolimits}
\def\proetale{{pro\text{-}\acute{e}tale}}
\def\etale{{\acute{e}tale}}
\def\QCoh{\mathit{QCoh}}
\def\Ker{\mathop{\mathrm{Ker}}}
\def\Im{\mathop{\mathrm{Im}}}
\def\Coker{\mathop{\mathrm{Coker}}}
\def\Coim{\mathop{\mathrm{Coim}}}

% Boxtimes
%
\DeclareMathSymbol{\boxtimes}{\mathbin}{AMSa}{"02}

%
% Macros for moduli stacks/spaces
%
\def\QCohstack{\mathcal{QC}\!\mathit{oh}}
\def\Cohstack{\mathcal{C}\!\mathit{oh}}
\def\Spacesstack{\mathcal{S}\!\mathit{paces}}
\def\Quotfunctor{\mathrm{Quot}}
\def\Hilbfunctor{\mathrm{Hilb}}
\def\Curvesstack{\mathcal{C}\!\mathit{urves}}
\def\Polarizedstack{\mathcal{P}\!\mathit{olarized}}
\def\Complexesstack{\mathcal{C}\!\mathit{omplexes}}
% \Pic is the operator that assigns to X its picard group, usage \Pic(X)
% \Picardstack_{X/B} denotes the Picard stack of X over B
% \Picardfunctor_{X/B} denotes the Picard functor of X over B
\def\Pic{\mathop{\mathrm{Pic}}\nolimits}
\def\Picardstack{\mathcal{P}\!\mathit{ic}}
\def\Picardfunctor{\mathrm{Pic}}
\def\Deformationcategory{\mathcal{D}\!\mathit{ef}}
\usepackage{cleveref}
\crefname{originallemma}{original lemma}{original lemmas}
\Crefname{originallemma}{Original lemma}{Original lemmas}
\makeatletter
\newlabel{lemma-cohomology-base-change}{{08EU}{source}{}{sourcepointer-1}{}}
\newlabel{lemma-cohomology-base-change@cref}{{[originallemma][0][]08EU}{[1][0][]source}{}{}{}}
\makeatother
\crefname{originallemma}{original lemma}{original lemmas}
\Crefname{originallemma}{Original lemma}{Original lemmas}
\makeatletter
\newlabel{lemma-quasi-coherence-pushforward-direct-sums}{{08DZ}{source}{}{sourcepointer-2}{}}
\newlabel{lemma-quasi-coherence-pushforward-direct-sums@cref}{{[originallemma][0][]08DZ}{[1][0][]source}{}{}{}}
\makeatother
\crefname{originallemma}{original lemma}{original lemmas}
\Crefname{originallemma}{Original lemma}{Original lemmas}
\makeatletter
\newlabel{cohomology-lemma-dual-perfect-complex}{{08DQ}{source}{}{sourcepointer-3}{}}
\newlabel{cohomology-lemma-dual-perfect-complex@cref}{{[originallemma][0][]08DQ}{[1][0][]source}{}{}{}}
\makeatother
\crefname{originallemma}{original lemma}{original lemmas}
\Crefname{originallemma}{Original lemma}{Original lemmas}
\makeatletter
\newlabel{limits-lemma-descend-finite-presentation}{{01ZM}{source}{}{sourcepointer-4}{}}
\newlabel{limits-lemma-descend-finite-presentation@cref}{{[originallemma][0][]01ZM}{[1][0][]source}{}{}{}}
\makeatother
\crefname{originallemma}{original lemma}{original lemmas}
\Crefname{originallemma}{Original lemma}{Original lemmas}
\makeatletter
\newlabel{limits-lemma-descend-flat-finite-presentation}{{04AI}{source}{}{sourcepointer-5}{}}
\newlabel{limits-lemma-descend-flat-finite-presentation@cref}{{[originallemma][0][]04AI}{[1][0][]source}{}{}{}}
\makeatother
\crefname{originallemma}{original lemma}{original lemmas}
\Crefname{originallemma}{Original lemma}{Original lemmas}
\makeatletter
\newlabel{lemma-compare-base-change}{{08IB}{source}{}{sourcepointer-6}{}}
\newlabel{lemma-compare-base-change@cref}{{[originallemma][0][]08IB}{[1][0][]source}{}{}{}}
\makeatother
\crefname{originallemma}{original lemma}{original lemmas}
\Crefname{originallemma}{Original lemma}{Original lemmas}
\makeatletter
\newlabel{lemma-countable}{{0CRQ}{source}{}{sourcepointer-7}{}}
\newlabel{lemma-countable@cref}{{[originallemma][0][]0CRQ}{[1][0][]source}{}{}{}}
\makeatother
\crefname{originallemma}{original lemma}{original lemmas}
\Crefname{originallemma}{Original lemma}{Original lemmas}
\makeatletter
\newlabel{lemma-countable-cohomology}{{0CRP}{source}{}{sourcepointer-8}{}}
\newlabel{lemma-countable-cohomology@cref}{{[originallemma][0][]0CRP}{[1][0][]source}{}{}{}}
\makeatother
\title{Sections over a flat finitely presented morphism as a homotopy colimit of perfect objects}
\author{The Stacks Project Authors}
\date{stacks-project:perfect.tex:local-snapshot-20261012; GNU FDL 1.2 or later}
\begin{document}
\maketitle
\setcounter{subsection}{1}
\noindent\textit{Editorial note.} All delivered mathematical text is the author's own, copied from the pinned Stacks Project chapter file perfect.tex; no statement or proof step was rewritten, reordered or removed. The counted span is the author environment and its own proof environment. References inside the span are handed to the assembler as navigation stubs that carry the canonical tag of the cited tag only; whether each of them is load-bearing is recorded on the stub itself.\par\medskip
\noindent\textit{Allowed standard prerequisites (editorial).}\begin{itemize}
\item For a flat morphism of finite presentation the cohomology of the derived pullback is computed by the derived tensor of the pushforward of the structure sheaf, functorially in K. used at proof lines 8629-8633. Converts RGamma(U, Lf\textasciicircum{}*K) into RGamma(X, Rf\_*O\_U tensor\textasciicircum{}L K), the input to the colimit argument.
\item The global sections functor on a quasi-compact separated scheme commutes with homotopy colimits because it commutes with direct sums, and tensoring with a fixed complex commutes with derived colimits because direct sums in the derived category are computed termwise. invoked at proof lines 8633-8641. Justifies passing the homotopy colimit through both sides of the desired identity.
\item Morphisms of finite presentation and flatness descend along the approximation of the base ring by finitely generated subrings, and the derived pullback of the pushforward recovers the original pushforward under base change. invoked at proof lines 8647-8659. Reduces to the case where U and X are of finite type over Z.
\item For a morphism of finite type over Z the pushforward of the structure sheaf is a homotopy colimit of perfect complexes as soon as the sheaves R\textasciicircum{}if\_*O\_U have countable sets of sections over affine opens, and the structure sheaf satisfies this countability. invoked at proof lines 8661-8670. Supplies the existence of the system of perfect objects E\_n in the finite type case.
\end{itemize}\medskip
\ifdefined\hypertarget\hypertarget{sourcepointer-1}{}\fi
\noindent\textit{Source reference.} Original lemma 08EU: sources/stacks-perfect/pinned-original-source.tex, source lines 5269--5269. Cited at proof lines 8629-8633 to rewrite RGamma(U, Lf\textasciicircum{}*K); the colimit reduction using it is inside the counted span.\par\medskip
\ifdefined\hypertarget\hypertarget{sourcepointer-2}{}\fi
\noindent\textit{Source reference.} Original lemma 08DZ: sources/stacks-perfect/pinned-original-source.tex, source lines 803--803. Cited at proof lines 8635-8637 for RGamma(X, -) commuting with direct sums; the conclusion that it commutes with homotopy colimits is drawn inside the counted span.\par\medskip
\ifdefined\hypertarget\hypertarget{sourcepointer-3}{}\fi
\noindent\textit{Source reference.} Original lemma 08DQ: sources/stacks-chow/tags.txt, source lines 1--1. Cited at proof lines 8643-8646 to obtain (1) from (2) by setting L\_n = E\_n\textasciicircum{}dual; the setting of the dual system is inside the counted span.\par\medskip
\ifdefined\hypertarget\hypertarget{sourcepointer-4}{}\fi
\noindent\textit{Source reference.} Original lemma 01ZM: sources/stacks-chow/tags.txt, source lines 1--1. Cited at proof lines 8649-8653 to descend U to X to Spec(A\_i); the use of the descended data is inside the counted span.\par\medskip
\ifdefined\hypertarget\hypertarget{sourcepointer-5}{}\fi
\noindent\textit{Source reference.} Original lemma 04AI: sources/stacks-chow/tags.txt, source lines 1--1. Cited at proof lines 8654-8656 to assume f\_i flat; the reduction to f\_i is inside the counted span.\par\medskip
\ifdefined\hypertarget\hypertarget{sourcepointer-6}{}\fi
\noindent\textit{Source reference.} Original lemma 08IB: sources/stacks-perfect/pinned-original-source.tex, source lines 5362--5362. Cited at proof lines 8657-8659 to identify Rf\_*O\_U with the pullback of Rf\_\{i,*\}O\_\{U\_i\}; the commutation of Lg\textasciicircum{}* with colimits and the resulting reduction are inside the counted span.\par\medskip
\ifdefined\hypertarget\hypertarget{sourcepointer-7}{}\fi
\noindent\textit{Source reference.} Original lemma 0CRQ: sources/stacks-perfect/pinned-original-source.tex, source lines 8557--8557. Cited at proof lines 8662-8665 to reduce the existence of E\_n to countability; the application and the countability conclusion are inside the counted span.\par\medskip
\ifdefined\hypertarget\hypertarget{sourcepointer-8}{}\fi
\noindent\textit{Source reference.} Original lemma 0CRP: sources/stacks-perfect/pinned-original-source.tex, source lines 8508--8508. Cited at proof lines 8666-8670 applied to the structure sheaf; the final assembly of the argument is inside the counted span.\par\medskip
\begin{lemma}
\label{lemma-computing-sections-as-colim}
Let $A$ be a ring. Let $X$ be a scheme of finite presentation over $A$.
Let $f : U \to X$ be a flat morphism of finite presentation. Then
\begin{enumerate}
\item there exists an inverse system of perfect objects $L_n$ of
$D(\mathcal{O}_X)$ such that
$$
R\Gamma(U, Lf^*K) = \text{hocolim}\ R\Hom_X(L_n, K)
$$
in $D(A)$ functorially in $K$ in $D_\QCoh(\mathcal{O}_X)$, and
\item there exists a system of perfect objects $E_n$ of
$D(\mathcal{O}_X)$ such that
$$
R\Gamma(U, Lf^*K) = \text{hocolim}\ R\Gamma(X, E_n \otimes^\mathbf{L} K)
$$
in $D(A)$ functorially in $K$ in $D_\QCoh(\mathcal{O}_X)$.
\end{enumerate}
\end{lemma}

\begin{proof}
By Lemma \ref{lemma-cohomology-base-change} we have
$$
R\Gamma(U, Lf^*K) = R\Gamma(X, Rf_*\mathcal{O}_U \otimes^\mathbf{L} K)
$$
functorially in $K$. Observe that $R\Gamma(X, -)$ commutes with
homotopy colimits because it commutes with direct sums by
Lemma \ref{lemma-quasi-coherence-pushforward-direct-sums}.
Similarly, $- \otimes^\mathbf{L} K$ commutes with derived colimits
because $- \otimes^\mathbf{L} K$ commutes with direct sums
(because direct sums in $D(\mathcal{O}_X)$
are given by direct sums of representing complexes).
Hence to prove (2) it suffices to write
$Rf_*\mathcal{O}_U = \text{hocolim} E_n$ for a system of
perfect objects $E_n$ of $D(\mathcal{O}_X)$. Once this is done
we obtain (1) by setting $L_n = E_n^\vee$, see
Cohomology, Lemma \ref{cohomology-lemma-dual-perfect-complex}.

\medskip\noindent
Write $A = \colim A_i$ with $A_i$ of finite type over $\mathbf{Z}$.
By Limits, Lemma \ref{limits-lemma-descend-finite-presentation}
we can find an $i$ and morphisms $U_i \to X_i \to \Spec(A_i)$
of finite presentation whose base change to $\Spec(A)$ recovers
$U \to X \to \Spec(A)$.
After increasing $i$ we may assume that $f_i : U_i \to X_i$ is
flat, see Limits, Lemma \ref{limits-lemma-descend-flat-finite-presentation}.
By Lemma \ref{lemma-compare-base-change}
the derived pullback of $Rf_{i, *}\mathcal{O}_{U_i}$
by $g : X \to X_i$ is equal to $Rf_*\mathcal{O}_U$.
Since $Lg^*$ commutes with derived colimits, it suffices
to prove what we want for $f_i$. Hence we may assume that
$U$ and $X$ are of finite type over $\mathbf{Z}$.

\medskip\noindent
Assume $f : U \to X$ is a morphism of schemes of finite type over $\mathbf{Z}$.
To finish the proof we will show that $Rf_*\mathcal{O}_U$ is a homotopy
colimit of perfect complexes. To see this we apply Lemma \ref{lemma-countable}.
Thus it suffices to show that $R^if_*\mathcal{O}_U$
has countable sets of sections over affine opens.
This follows from Lemma \ref{lemma-countable-cohomology}
applied to the structure sheaf.
\end{proof}

\end{document}
